\section{Experimental Setup in Detail}
\label{sup:setup}

\subsection{Data}
Candidates were checked on their official pages on 2026-11-14 (no data downloaded). The setting needs
raw aerial images, a pose per image (exterior orientation, GNSS/IMU or registration) and a reference
surface that does not come from photogrammetry, since a photogrammetric reference would share the
errors under test.

\runin{UseGeo (main real set).} UAV nadir images from about 80~m, about 2~cm ground sampling distance,
829 images in three sets of eight strips each (224, 328 and 277 images; 80\% forward and 60\% side
overlap), exterior orientation and interior parameters per image, UAV LiDAR at about 50~points/m$^2$ and
per-image LiDAR depth maps \cite{data_usegeo}. Licence CC BY-NC-SA 4.0, as stated in the README of its
GitHub repository (the repository holds no separate licence file); download links per set, size not
stated. The LiDAR, not the also provided photogrammetric point cloud, is the reference.

\runin{H3D (second real set, pending).} UAV LiDAR at about 800~points/m$^2$ and textured meshes, in several
epochs, with oblique imagery at 1.5--3~cm used for the meshes \cite{data_h3d}. The download link is given
only after registration, the data are provided as is, and the pages state no licence and do not say
that the images or their orientations are released. \TODO{Ask the maintainers before relying on H3D;
treated as non-commercial research use.}

\runin{MatrixCity (synthetic).} A synthetic city rendered in a game engine, with aerial and street views
\cite{data_matrixcity}. The Hugging Face dataset card states CC BY-NC 4.0; the files include rendered depth,
normals and poses (\texttt{transforms.json}). It is used for controlled fault injection and the pose check
with exact truth. One aerial test block alone is about 15.8~GB.

\runin{UrbanScene3D (optional, one scene at most).} Real and synthetic drone scenes, some with LiDAR
\cite{data_urbanscene3d}, under custom non-commercial terms that also forbid disseminating the data or
altered versions; 1.43~TB in total. At most one real scene is used, for comparability with large-scene
splatting work. \TODO{Choose a scene whose LiDAR and poses are confirmed.}

\runin{ETH3D (sanity check).} Multi-view images with laser-scan ground truth and registered poses
\cite{data_eth3d}, used as given. \TODO{Fix the scene subset and check the licence.}

\runin{Considered and not used.} ISPRS Vaihingen (crewed nadir frames with exterior orientation and
airborne LiDAR; research use only, no redistribution, 16~GB) is the fallback real set. The ISPRS/EuroSDR
Dortmund benchmark has oblique and UAV images with GNSS/INS and laser scans, but its oblique blocks are
125--350~GB and the scan type is not stated. Mill-19 has no geometric reference and could serve only as a
rendering check. ScanNet++ was dropped: it is indoor, and the aerial sets cover the specular and
textureless failures that motivated it. All datasets are licensed for non-commercial use only, and no
large download is made without the author's approval.

\subsection{Depth Model}
Depth Pro \cite{est_depthpro} is the default: it predicts metric depth without intrinsics, one map per
image, and it returns no confidence, so its $\sigma(\cdot)$ is fitted (H1). Per their public
documentation (checked 2026-11-14), Metric3D v2 returns a depth confidence for its ViT models and
UniDepthV2 lists a confidence output; they supply the confidence baseline. MoGe returns a validity
mask rather than a confidence, and Depth Anything V2 and Marigold document no uncertainty output.

\subsection{Reconstruction and Gate}
\runin{Poses.} The input is posed images. On aerial data the trajectory is a uniform cubic B-spline in time,
with one control point per second of flight (fixed on development data), fitted to the exterior
orientation or GNSS/IMU with a second-difference smoothness prior; rotations use the cumulative form on
$SO(3)$. The control points are then refined by differentiable pose-graph optimisation whose factors are
the GNSS/IMU prior, the smoothness prior and render-and-compare depth terms, and optionally photometric
render-and-compare of the appearance block; no feature correspondences are used. Before a view is
tested, its own render-and-compare term is subtracted from the control-point information, and its pose
covariance $\Sigma_q$ follows from the spline's basis row at its time stamp. The ablation starts from poses
perturbed by a fixed amount (fixed on development data). On ETH3D the registered poses are used as given.

\runin{Blocks.} Blocks are $B$ consecutive frames of the same flight strip, centred on the tested view and
shifted inwards at the strip ends; strips are taken from the flight plan or split at turns of the spline.
$B$ is chosen from $\{1,3,5,9\}$ on development data against the measured decay of the cross-view error
correlation with frame separation (H2).

\runin{Train, freeze, fuse, test.} Depth maps are computed once and cached, so the depth model is not
resident during training. Surfels are 2DGS planar Gaussians rendered with gsplat into expected depth,
alpha weights and accumulated opacity. During a warm-up of a few thousand iterations (fixed on
development scenes) the centres, rotations, scales and opacities are fitted on the depth loss over all
views, $(a_i,b_i)$ in closed form and $\beta$ by least squares; densification needs at least two
supporting views and pruning removes surfels below a minimum opacity. Opacity receives depth gradients
only. The opacity is then frozen, densification and pruning stop, and the geometry and the poses are fitted
to convergence while the colours train on the frozen geometry. The Laplace information \eqref{M-eq:info}
is accumulated at the MAP, and each view is tested by Algorithm~\ref{M-alg:gate}.

\runin{Outer loop.} Rounds of train, freeze, fuse, test as in Algorithm~\ref{M-alg:gate}. Each round warm-starts
from the previous fit, unfreezes opacity and re-enables densification and pruning, and trains only on
the active views. Defaults: one view dropped per round, no readmission, $R_{\max}=10$, Holm step-down over
$M=NR_{\max}$ tests. Ablations: loop off ($R_{\max}=1$), $R_{\max}=3$, dropping all rejected views per
round, down-weighting a flagged view to a tenth instead of dropping it, and readmission (a dropped view
re-tested each round and readmitted once if not rejected).

\runin{Parameters.} The support threshold $\lambda_{\min}$, the opacity share $\rho=\tfrac12$, the
opacity threshold $\tau_W$, the prior $\Lambda_0$, the warm-up length, the candidate regions of the scan
(appearance segments), the temperature range for the scene-level alarm and the 5\% per-view false alarm
are fixed on development scenes. \TODO{Fix them; report GPU model, wall-clock and memory on the 8~GB card.}

\runin{Pose perturbation check (H7).} After the gate has run on a scene, the poses of chosen views are
perturbed: isolated clean views, a block of consecutive clean views, and the flagged view. Each
perturbation is a rotation of a known angle about a random axis, from $0.05^{\circ}$ to $2^{\circ}$, and a
translation of a known length in a random direction, from 1 to 50 ground sampling distances. (a)~The
gate is run before any refinement, with each perturbed view's own depth excluded from its pose as
always; (b)~the spline, pose-graph and render-and-compare refinement is run, and the pose error to the
reference (exact on MatrixCity, the given exterior orientation on UseGeo) is measured as an angle and a
distance. Detection rate and recovery error are reported against perturbation size.

\subsection{Depth-Model Altitude Comparison}
Depth Pro, Metric3D v2 and UniDepthV2 are run as released on nadir views from survey altitude and on
near-ground views (ETH3D, and low oblique frames where available). For each, the fraction of views with a
regional failure, the failure area, its magnitude after alignment and the fitted noise level are
compared between the two regimes (H1, Table~\ref{tab:altitude}).

\subsection{Faults}
\runin{Natural.} A view region is a natural failure if its aligned error against the laser depth
exceeds $k\sigma$ on a connected region of at least a minimum area; regions are classed from the image
and the reference as specular (glass, water), textureless (roofs, roads, walls), or other. $k$ and the minimum area
are fixed on development scenes.
\runin{Injected.} Offset, plane replacement, blurred copy and mirrored depth (the depth of a reflected
scene) on 2, 5, 10, 20 or 40\% of a view at 2, 4, 8 or 16$\sigma$, placed on specular or textureless
appearance masks or on random connected superpixels, either independently per view or repeated with
the same error across views. Whole-view affine errors are a control: they are absorbed and should be
neither detected nor harmful.

\subsection{Baselines}
\emph{Robust splatting:} SpotLessSplats as released with depth supervision added \cite{gs_spotless};
the RobustNeRF trimmed robust loss applied to splatting \cite{gs_robustnerf}; Huber and
Geman--McClure depth losses \cite{rob_blackrangarajan,rob_barron}. \emph{Cheap view tests:} the RMS of
the per-view aligned depth residual against the other views' render, without whitening; and a view
score from the confidence of Metric3D v2 or UniDepthV2. For H3, the masks of the robust methods are
also turned into view scores. \emph{Bounds:} no gate, and an oracle gate that drops exactly the faulty
views. \emph{Gate ablations:} a joint surface (no leave-view-out); keeping surfels that only the tested
view supports; constant-$\sigma$ residuals; information deflated by $n_{\mathrm{eff}}/n$; jackknife
inflation of the rendered variance; no scene temperature; no shared distance bias; a fold-wise opacity
warm-up ($K=5$ warm-ups, each without one fold); the outer-loop ablations above; single-view leave-out ($B=1$) and the block-size sweep; perturbed starting
poses; a FisherRF-style
diagonal Fisher \cite{unc_fisherrf} and a Bayes' Rays-style uncertainty \cite{unc_bayesrays} in place
of the rendered blocks; and energy or region scan alone.

\subsection{Metrics and Protocol}
Detection power at 5\% per-view false alarm, false-alarm, untestable and availability rates and the
elimination-curve AUC; held-out laser-depth AbsRel and surface accuracy, completeness and F-score; the
cross-view error correlation by appearance class; raw and temperature-scaled $T_j/m_j$ on clean views, the fitted $\hat\gamma$ per scene and $\omega$.
Scene-level bootstrap intervals and paired Wilcoxon tests, five seeds. Hypotheses, thresholds and the
dev/test split are fixed in a pre-registration file and tagged before any test scene is processed.
